% --- Archon physics formalization source begin ---
% archon:physics
% archon:covers PhyXMiniProblems/problem_phyx_mini_0056.lean
% archon:source-report reports/phyx_mini/problem_phyx_mini_0056.source.json

\chapter{Physics problem phyx\_mini\_0056}
\label{ch:PhyXMiniProblems_problem_phyx_mini_0056}

\paragraph{Problem source.}
\#\# Question

What is the focal length of the glass meniscus lens shown?

\#\# Answer choices

- A: A: 14cm

- B: B: 86cm

- C: C: 80cm

- D: D: 74cm

\#\# Figure caption (auxiliary; use the image as primary evidence)

The image depicts a concave lens with light gray and blue shading, illustrating its curvature. There are two radii extending from the center of the curvature to the edges of the lens. 

- **Radii:**
  - \textbackslash{}( R\_1 = 40 \textbackslash{}, \textbackslash{}text\{cm\} \textbackslash{}) on the left side.
  - \textbackslash{}( R\_2 = 20 \textbackslash{}, \textbackslash{}text\{cm\} \textbackslash{}) on the right side.

- **Refractive Index:**
  - The lens has a refractive index \textbackslash{}( n = 1.50 \textbackslash{}).

These elements are labeled accordingly. The scene reflects a simplified diagram of optical properties related to the lens.

\#\# Dataset metadata

Geometrical Optics; Physical Model Grounding Reasoning, Multi-Formula Reasoning

\paragraph{Recorded answer/context.}
C: C: 80cm

\paragraph{Figure/image path.}
/root/proposal\_for\_physic/science-mango/phyx\_mini\_run/phyx\_data/test\_image/56.png

\paragraph{Formalization target.}
create a compiling Lean file with sorry bodies at `PhyXMiniProblems/problem\_phyx\_mini\_0056.lean`. The Lean declarations must preserve the physical quantities, dimensions or dimensional roles, figure labels, governing-law hypotheses, and final relation expressed by this problem.
Use Mathlib/Physlib names found through LeanExplore where available. If a physics API is missing, introduce faithful local abstractions rather than scalar placeholder aliases.

\begin{theorem}[Physics formalization target]
\label{thm:physics:phyx_mini_0056:target}
\lean{PhyXMiniProblems.ProblemPhyXMini0056.focal_length_of_glass_meniscus_lens}
\uses{def:physics:phyx-mini-0056:phyxminiproblems-problemphyxmini0056-signedlengthincentimeters, def:physics:phyx-mini-0056:phyxminiproblems-problemphyxmini0056-glassmeniscuslens, def:physics:phyx-mini-0056:phyxminiproblems-problemphyxmini0056-matchesfigurereadout, def:physics:phyx-mini-0056:phyxminiproblems-problemphyxmini0056-modelsambientair, def:physics:phyx-mini-0056:phyxminiproblems-problemphyxmini0056-hasphysicalopticalparameters, def:physics:phyx-mini-0056:phyxminiproblems-problemphyxmini0056-obeysthinlensmakerequation, def:physics:phyx-mini-0056:phyxminiproblems-problemphyxmini0056-matchesanswerchoice}
The assigned autoformalize agent should translate this physics problem into Lean declarations in the covered file, with theorem and lemma proof bodies written as `by sorry`.
\end{theorem}
\begin{proof}
This is an autoformalization task, not a proof task. Produce statements that can later be proved by the physics prover without weakening the physical model.
\end{proof}

% --- Archon declaration topology begin ---
\section{Declaration topology}
\begin{definition}[Length Magnitude]
  \label{def:physics:phyx-mini-0056:phyxminiproblems-problemphyxmini0056-lengthmagnitude}
  \lean{PhyXMiniProblems.ProblemPhyXMini0056.LengthMagnitude}
  A nonnegative physical length magnitude, independent of readout units
\end{definition}
\begin{proof}
  This is the declared model definition of Length Magnitude.
\end{proof}
\begin{definition}[Signed Length Quantity]
  \label{def:physics:phyx-mini-0056:phyxminiproblems-problemphyxmini0056-signedlengthquantity}
  \lean{PhyXMiniProblems.ProblemPhyXMini0056.SignedLengthQuantity}
  A signed physical length, independent of readout units
\end{definition}
\begin{proof}
  This is the declared model definition of Signed Length Quantity.
\end{proof}
\begin{definition}[centimeter Unit Choices]
  \label{def:physics:phyx-mini-0056:phyxminiproblems-problemphyxmini0056-centimeterunitchoices}
  \lean{PhyXMiniProblems.ProblemPhyXMini0056.centimeterUnitChoices}
  Unit choices whose length component is centimeters and whose others are SI
\end{definition}
\begin{proof}
  This is the declared model definition of centimeter Unit Choices.
\end{proof}
\begin{definition}[length Magnitude Readout]
  \label{def:physics:phyx-mini-0056:phyxminiproblems-problemphyxmini0056-lengthmagnitudereadout}
  \lean{PhyXMiniProblems.ProblemPhyXMini0056.lengthMagnitudeReadout}
  \uses{def:physics:phyx-mini-0056:phyxminiproblems-problemphyxmini0056-lengthmagnitude}
  The scalar readout of a nonnegative physical length in a chosen unit system
\end{definition}
\begin{proof}
  This is the declared model definition of length Magnitude Readout.
\end{proof}
\begin{definition}[signed Length Readout]
  \label{def:physics:phyx-mini-0056:phyxminiproblems-problemphyxmini0056-signedlengthreadout}
  \lean{PhyXMiniProblems.ProblemPhyXMini0056.signedLengthReadout}
  \uses{def:physics:phyx-mini-0056:phyxminiproblems-problemphyxmini0056-signedlengthquantity}
  The scalar readout of a signed physical length in a chosen unit system
\end{definition}
\begin{proof}
  This is the declared model definition of signed Length Readout.
\end{proof}
\begin{definition}[length Magnitude In Centimeters]
  \label{def:physics:phyx-mini-0056:phyxminiproblems-problemphyxmini0056-lengthmagnitudeincentimeters}
  \lean{PhyXMiniProblems.ProblemPhyXMini0056.lengthMagnitudeInCentimeters}
  \uses{def:physics:phyx-mini-0056:phyxminiproblems-problemphyxmini0056-lengthmagnitude, def:physics:phyx-mini-0056:phyxminiproblems-problemphyxmini0056-centimeterunitchoices, def:physics:phyx-mini-0056:phyxminiproblems-problemphyxmini0056-lengthmagnitudereadout}
  The scalar centimeter readout of a nonnegative physical length
\end{definition}
\begin{proof}
  This is the declared model definition of length Magnitude In Centimeters.
\end{proof}
\begin{definition}[signed Length In Centimeters]
  \label{def:physics:phyx-mini-0056:phyxminiproblems-problemphyxmini0056-signedlengthincentimeters}
  \lean{PhyXMiniProblems.ProblemPhyXMini0056.signedLengthInCentimeters}
  \uses{def:physics:phyx-mini-0056:phyxminiproblems-problemphyxmini0056-signedlengthquantity, def:physics:phyx-mini-0056:phyxminiproblems-problemphyxmini0056-centimeterunitchoices, def:physics:phyx-mini-0056:phyxminiproblems-problemphyxmini0056-signedlengthreadout}
  The scalar centimeter readout of a signed physical length
\end{definition}
\begin{proof}
  This is the declared model definition of signed Length In Centimeters.
\end{proof}
\begin{definition}[Axial Side]
  \label{def:physics:phyx-mini-0056:phyxminiproblems-problemphyxmini0056-axialside}
  \lean{PhyXMiniProblems.ProblemPhyXMini0056.AxialSide}
  The side of a surface vertex on which its center of curvature lies
\end{definition}
\begin{proof}
  This is the declared model definition of Axial Side.
\end{proof}
\begin{definition}[Propagation Direction]
  \label{def:physics:phyx-mini-0056:phyxminiproblems-problemphyxmini0056-propagationdirection}
  \lean{PhyXMiniProblems.ProblemPhyXMini0056.PropagationDirection}
  Propagation directions along the horizontal optical axis
\end{definition}
\begin{proof}
  This is the declared model definition of Propagation Direction.
\end{proof}
\begin{definition}[Surface Radius Label]
  \label{def:physics:phyx-mini-0056:phyxminiproblems-problemphyxmini0056-surfaceradiuslabel}
  \lean{PhyXMiniProblems.ProblemPhyXMini0056.SurfaceRadiusLabel}
  The two radius labels printed in the lens diagram
\end{definition}
\begin{proof}
  This is the declared model definition of Surface Radius Label.
\end{proof}
\begin{definition}[Spherical Surface Shape]
  \label{def:physics:phyx-mini-0056:phyxminiproblems-problemphyxmini0056-sphericalsurfaceshape}
  \lean{PhyXMiniProblems.ProblemPhyXMini0056.SphericalSurfaceShape}
  Surface shapes used to describe the two faces of the pictured meniscus
\end{definition}
\begin{proof}
  This is the declared model definition of Spherical Surface Shape.
\end{proof}
\begin{definition}[Lens Profile]
  \label{def:physics:phyx-mini-0056:phyxminiproblems-problemphyxmini0056-lensprofile}
  \lean{PhyXMiniProblems.ProblemPhyXMini0056.LensProfile}
  Profiles used to identify the pictured glass element
\end{definition}
\begin{proof}
  This is the declared model definition of Lens Profile.
\end{proof}
\begin{definition}[Optical Medium]
  \label{def:physics:phyx-mini-0056:phyxminiproblems-problemphyxmini0056-opticalmedium}
  \lean{PhyXMiniProblems.ProblemPhyXMini0056.OpticalMedium}
  Optical media occurring in the thin-lens model
\end{definition}
\begin{proof}
  This is the declared model definition of Optical Medium.
\end{proof}
\begin{definition}[Spherical Lens Surface]
  \label{def:physics:phyx-mini-0056:phyxminiproblems-problemphyxmini0056-sphericallenssurface}
  \lean{PhyXMiniProblems.ProblemPhyXMini0056.SphericalLensSurface}
  \uses{def:physics:phyx-mini-0056:phyxminiproblems-problemphyxmini0056-lengthmagnitude, def:physics:phyx-mini-0056:phyxminiproblems-problemphyxmini0056-axialside, def:physics:phyx-mini-0056:phyxminiproblems-problemphyxmini0056-surfaceradiuslabel, def:physics:phyx-mini-0056:phyxminiproblems-problemphyxmini0056-sphericalsurfaceshape}
  A spherical refracting surface retains its unsigned physical radius magnitude. The sign used by the optical convention is recovered from the propagation direction and the center-of-curvature side
\end{definition}
\begin{proof}
  This is the declared model definition of Spherical Lens Surface.
\end{proof}
\begin{definition}[Glass Meniscus Lens]
  \label{def:physics:phyx-mini-0056:phyxminiproblems-problemphyxmini0056-glassmeniscuslens}
  \lean{PhyXMiniProblems.ProblemPhyXMini0056.GlassMeniscusLens}
  \uses{def:physics:phyx-mini-0056:phyxminiproblems-problemphyxmini0056-signedlengthquantity, def:physics:phyx-mini-0056:phyxminiproblems-problemphyxmini0056-propagationdirection, def:physics:phyx-mini-0056:phyxminiproblems-problemphyxmini0056-surfaceradiuslabel, def:physics:phyx-mini-0056:phyxminiproblems-problemphyxmini0056-lensprofile, def:physics:phyx-mini-0056:phyxminiproblems-problemphyxmini0056-opticalmedium, def:physics:phyx-mini-0056:phyxminiproblems-problemphyxmini0056-sphericallenssurface}
  The glass meniscus lens, its two labeled surfaces, material indices, incident direction, and signed focal length
\end{definition}
\begin{proof}
  This is the declared model definition of Glass Meniscus Lens.
\end{proof}
\begin{definition}[signed Radius Readout]
  \label{def:physics:phyx-mini-0056:phyxminiproblems-problemphyxmini0056-signedradiusreadout}
  \lean{PhyXMiniProblems.ProblemPhyXMini0056.signedRadiusReadout}
  \uses{def:physics:phyx-mini-0056:phyxminiproblems-problemphyxmini0056-lengthmagnitudereadout, def:physics:phyx-mini-0056:phyxminiproblems-problemphyxmini0056-propagationdirection, def:physics:phyx-mini-0056:phyxminiproblems-problemphyxmini0056-sphericallenssurface}
  The signed scalar radius readout in the coordinate direction of propagation. A center on the incident side is negative and one on the outgoing side is positive
\end{definition}
\begin{proof}
  This is the declared model definition of signed Radius Readout.
\end{proof}
\begin{definition}[signed Radius In Centimeters]
  \label{def:physics:phyx-mini-0056:phyxminiproblems-problemphyxmini0056-signedradiusincentimeters}
  \lean{PhyXMiniProblems.ProblemPhyXMini0056.signedRadiusInCentimeters}
  \uses{def:physics:phyx-mini-0056:phyxminiproblems-problemphyxmini0056-centimeterunitchoices, def:physics:phyx-mini-0056:phyxminiproblems-problemphyxmini0056-propagationdirection, def:physics:phyx-mini-0056:phyxminiproblems-problemphyxmini0056-sphericallenssurface, def:physics:phyx-mini-0056:phyxminiproblems-problemphyxmini0056-signedradiusreadout}
  The signed radius readout of a surface in centimeters
\end{definition}
\begin{proof}
  This is the declared model definition of signed Radius In Centimeters.
\end{proof}
\begin{definition}[Matches Figure Readout]
  \label{def:physics:phyx-mini-0056:phyxminiproblems-problemphyxmini0056-matchesfigurereadout}
  \lean{PhyXMiniProblems.ProblemPhyXMini0056.MatchesFigureReadout}
  \uses{def:physics:phyx-mini-0056:phyxminiproblems-problemphyxmini0056-lengthmagnitudeincentimeters, def:physics:phyx-mini-0056:phyxminiproblems-problemphyxmini0056-glassmeniscuslens}
  Data read from the primary image: lens profile and surface geometry, the two radius labels and magnitudes, left-to-right propagation convention, and the dimensionless glass refractive index 1.50. No focal-length value occurs here
\end{definition}
\begin{proof}
  This is the declared model definition of Matches Figure Readout.
\end{proof}
\begin{definition}[Models Ambient Air]
  \label{def:physics:phyx-mini-0056:phyxminiproblems-problemphyxmini0056-modelsambientair}
  \lean{PhyXMiniProblems.ProblemPhyXMini0056.ModelsAmbientAir}
  \uses{def:physics:phyx-mini-0056:phyxminiproblems-problemphyxmini0056-glassmeniscuslens}
  The usual environmental idealization that the pictured lens is surrounded by air of dimensionless refractive index one
\end{definition}
\begin{proof}
  This is the declared model definition of Models Ambient Air.
\end{proof}
\begin{definition}[Has Physical Optical Parameters]
  \label{def:physics:phyx-mini-0056:phyxminiproblems-problemphyxmini0056-hasphysicalopticalparameters}
  \lean{PhyXMiniProblems.ProblemPhyXMini0056.HasPhysicalOpticalParameters}
  \uses{def:physics:phyx-mini-0056:phyxminiproblems-problemphyxmini0056-lengthmagnitudeincentimeters, def:physics:phyx-mini-0056:phyxminiproblems-problemphyxmini0056-glassmeniscuslens}
  Positivity and optical-density conditions for ordinary glass in air. These conditions contain no value for the requested focal length
\end{definition}
\begin{proof}
  This is the declared model definition of Has Physical Optical Parameters.
\end{proof}
\begin{definition}[Obeys Thin Lensmaker Equation]
  \label{def:physics:phyx-mini-0056:phyxminiproblems-problemphyxmini0056-obeysthinlensmakerequation}
  \lean{PhyXMiniProblems.ProblemPhyXMini0056.ObeysThinLensmakerEquation}
  \uses{def:physics:phyx-mini-0056:phyxminiproblems-problemphyxmini0056-signedlengthreadout, def:physics:phyx-mini-0056:phyxminiproblems-problemphyxmini0056-glassmeniscuslens, def:physics:phyx-mini-0056:phyxminiproblems-problemphyxmini0056-signedradiusreadout}
  The thin-lens lensmaker equation in a uniform surrounding medium, 1 / f = (n\_lens / n\_medium - 1) * (1 / R₁ - 1 / R₂). It is stated for every unit choice, so focal length and radii remain physical lengths rather than being identified with bare centimeter scalars. This is a generic governing-law hypothesis and contains no numerical focal-length answer
\end{definition}
\begin{proof}
  This is the declared model definition of Obeys Thin Lensmaker Equation.
\end{proof}
\begin{definition}[Answer Choice]
  \label{def:physics:phyx-mini-0056:phyxminiproblems-problemphyxmini0056-answerchoice}
  \lean{PhyXMiniProblems.ProblemPhyXMini0056.AnswerChoice}
  Labels of the four multiple-choice answers displayed with the problem
\end{definition}
\begin{proof}
  This is the declared model definition of Answer Choice.
\end{proof}
\begin{definition}[focal Length In Centimeters]
  \label{def:physics:phyx-mini-0056:phyxminiproblems-problemphyxmini0056-answerchoice-focallengthincentimeters}
  \lean{PhyXMiniProblems.ProblemPhyXMini0056.AnswerChoice.focalLengthInCentimeters}
  \uses{def:physics:phyx-mini-0056:phyxminiproblems-problemphyxmini0056-answerchoice}
  The centimeter focal-length readout printed beside each answer choice
\end{definition}
\begin{proof}
  This is the declared model definition of focal Length In Centimeters.
\end{proof}
\begin{definition}[Matches Answer Choice]
  \label{def:physics:phyx-mini-0056:phyxminiproblems-problemphyxmini0056-matchesanswerchoice}
  \lean{PhyXMiniProblems.ProblemPhyXMini0056.MatchesAnswerChoice}
  \uses{def:physics:phyx-mini-0056:phyxminiproblems-problemphyxmini0056-signedlengthincentimeters, def:physics:phyx-mini-0056:phyxminiproblems-problemphyxmini0056-glassmeniscuslens, def:physics:phyx-mini-0056:phyxminiproblems-problemphyxmini0056-answerchoice}
  A lens's centimeter focal-length readout agrees with a displayed choice
\end{definition}
\begin{proof}
  This is the declared model definition of Matches Answer Choice.
\end{proof}
\begin{lemma}[signed radius readouts of matches figure]
  \label{lem:physics:phyx-mini-0056:phyxminiproblems-problemphyxmini0056-signed-radius-readouts-of-matches-figure}
  \lean{PhyXMiniProblems.ProblemPhyXMini0056.signed_radius_readouts_of_matches_figure}
  \uses{def:physics:phyx-mini-0056:phyxminiproblems-problemphyxmini0056-glassmeniscuslens, def:physics:phyx-mini-0056:phyxminiproblems-problemphyxmini0056-signedradiusincentimeters, def:physics:phyx-mini-0056:phyxminiproblems-problemphyxmini0056-matchesfigurereadout}
  The pictured center locations turn the two nonnegative radius magnitudes into the signed lensmaker readouts R₁ = -40 cm and R₂ = -20 cm
\end{lemma}
\begin{proof}
  The conclusion follows from the governing relations and figure readouts named in the statement of signed radius readouts of matches figure.
\end{proof}
\begin{lemma}[reciprocal focal length centimeters]
  \label{lem:physics:phyx-mini-0056:phyxminiproblems-problemphyxmini0056-reciprocal-focal-length-centimeters}
  \lean{PhyXMiniProblems.ProblemPhyXMini0056.reciprocal_focal_length_centimeters}
  \uses{def:physics:phyx-mini-0056:phyxminiproblems-problemphyxmini0056-signedlengthincentimeters, def:physics:phyx-mini-0056:phyxminiproblems-problemphyxmini0056-glassmeniscuslens, def:physics:phyx-mini-0056:phyxminiproblems-problemphyxmini0056-matchesfigurereadout, def:physics:phyx-mini-0056:phyxminiproblems-problemphyxmini0056-modelsambientair, def:physics:phyx-mini-0056:phyxminiproblems-problemphyxmini0056-obeysthinlensmakerequation}
  Specializing the general lensmaker law to the image data and ambient air gives reciprocal focal length 1/80 cm⁻¹
\end{lemma}
\begin{proof}
  The conclusion follows from the governing relations and figure readouts named in the statement of reciprocal focal length centimeters.
\end{proof}
% --- Archon declaration topology end ---
% --- Archon physics formalization source end ---
